% PKG-4 "Tokens first, dark-ground video" -- one-page print specimen (design-packages-B, 2026-10-09)
% Fonts: STIX Two Text + STIX Two Math (OFL), Inter (OFL), Iosevka (OFL). Colours from ../tokens/pkg4-tokens.sty (generated).
\documentclass[fontsize=10.5pt,parskip=false]{scrartcl}
\usepackage[a4paper,inner=22mm,textwidth=108mm,marginparsep=6mm,marginparwidth=50mm,
  top=22mm,bottom=26mm,twoside=false]{geometry}
\usepackage{fontspec}
\usepackage{amsmath}
\usepackage[math-style=ISO,bold-style=ISO]{unicode-math}
\setmainfont{STIXTwoText}[Extension=.otf,UprightFont=*-Regular,ItalicFont=*-Italic,
  BoldFont=*-SemiBold,BoldItalicFont=*-SemiBoldItalic,Numbers=Proportional]
\setmathfont{STIXTwoMath-Regular.otf}
\setsansfont{Inter}[UprightFont=* Regular,BoldFont=* SemiBold,ItalicFont=* Italic,Scale=MatchLowercase,Numbers=Tabular]
\setmonofont{Iosevka}[Scale=MatchLowercase]
\usepackage[ngerman]{babel}
\usepackage[autostyle]{csquotes}
\usepackage[protrusion=true,expansion=false]{microtype}
\usepackage[locale=DE]{siunitx}
\usepackage{xcolor}
% pkg4-tokens.sty -- GENERATED by build_tokens.py from pkg4.tokens.json. Do not edit.
\definecolor{tpc-paper}{HTML}{FFFFFF}
\definecolor{tpcvid-paper}{HTML}{090F14}
\definecolor{tpc-paper-alt}{HTML}{F0F4F7}
\definecolor{tpcvid-paper-alt}{HTML}{131B20}
\definecolor{tpc-ink}{HTML}{111921}
\definecolor{tpcvid-ink}{HTML}{EBEFF2}
\definecolor{tpc-ink-muted}{HTML}{545C63}
\definecolor{tpcvid-ink-muted}{HTML}{ABB2B6}
\definecolor{tpc-rule}{HTML}{BFC5CA}
\definecolor{tpcvid-rule}{HTML}{373E44}
\definecolor{tpc-accent}{HTML}{2E5B75}
\definecolor{tpcvid-accent}{HTML}{73CCEA}
\definecolor{tpc-link}{HTML}{066082}
\definecolor{tpcvid-link}{HTML}{71D4EF}
\definecolor{tpc-state}{HTML}{003F60}
\definecolor{tpcvid-state}{HTML}{63D4FE}
\definecolor{tpc-dynamics}{HTML}{B12B35}
\definecolor{tpcvid-dynamics}{HTML}{FD736F}
\definecolor{tpc-effect}{HTML}{C68206}
\definecolor{tpcvid-effect}{HTML}{FFD167}
\definecolor{tpc-environment}{HTML}{CCCAC5}
\definecolor{tpcvid-environment}{HTML}{43494D}

\usepackage{tikz}
\usetikzlibrary{patterns,arrows.meta,calc}
\usepackage[most]{tcolorbox}
\usepackage{marginnote}
\color{tpc-ink}
\linespread{1.10}  % 10.5 on about 13.9
\setlength{\parindent}{1.2em}\setlength{\parskip}{0pt}
\newcommand{\dd}{\mathrm{d}}\newcommand{\eu}{\mathrm{e}}\newcommand{\iu}{\mathrm{i}}
\newcommand{\tr}{\operatorname{tr}}
\newcommand{\St}[1]{{\color{tpc-state}#1}}      % state role
\newcommand{\Ef}[1]{{\color{tpc-effect!80!black}#1}}  % effect role (text-weight variant)
% margin notes: Inter 8/11, ragged, ink-muted
\newcommand{\side}[1]{\marginpar{\raggedright\sffamily\fontsize{7.8}{10.6}\selectfont\color{tpc-ink-muted}#1}}
% theorem: accent rule left, paper-alt fill, sans run-in label
\newtcolorbox{postulat}[1]{enhanced,breakable,colback=tpc-paper-alt,colframe=tpc-paper-alt,
  boxrule=0pt,arc=0pt,left=4mm,right=3mm,top=2mm,bottom=2mm,
  borderline west={2pt}{0pt}{tpc-accent},
  fonttitle=\sffamily\bfseries\footnotesize,coltitle=tpc-accent,
  title={#1},attach title to upper={\quad},colbacktitle=tpc-paper-alt}
% Vertiefung (depth) block, after layered-docs/REPORT.md: outer margin rule + run-in label + tag in the margin
\newtcolorbox{zoomin}[2]{enhanced,breakable,colback=tpc-paper,colframe=tpc-paper,boxrule=0pt,arc=0pt,
  left=0pt,right=0pt,top=1mm,bottom=1mm,fontupper=\small,
  overlay={\draw[tpc-accent,line width=1.2pt]([xshift=3mm]frame.north east)--([xshift=3mm]frame.south east);
    \node[anchor=north west,inner sep=0pt,font=\sffamily\scriptsize\bfseries,text=tpc-accent,align=left]
      at ([xshift=6mm]frame.north east){VERTIEFUNG #1\\[-1pt]\mdseries\color{tpc-ink-muted}#2};}}
\pagestyle{empty}
\begin{document}
\noindent{\sffamily\bfseries\fontsize{8}{10}\selectfont\color{tpc-accent}KAPITEL 6 \enspace POSTULATE}\\[2mm]
{\fontsize{19}{23}\selectfont Die Bornsche Regel}\\[1mm]
{\sffamily\fontsize{9}{12}\selectfont\color{tpc-ink-muted}Wie ein Zustand und ein Detektor eine Wahrscheinlichkeit ergeben}
\vspace{5mm}

\noindent Eine Präparation liefert einen \emph{Zustand}, ein Detektor einen \emph{Effekt}. Die Theorie
sagt nicht voraus, \enquote{was passiert}, sondern mit welcher Wahrscheinlichkeit der Detektor
anspricht. Im spurklassigen Fall ist der Zustand ein Dichteoperator \(\St{\rho}\ge 0\) mit
\(\tr\St{\rho}=1\), der Effekt ein Operator \(0\le\Ef{E}\le\mathbb{1}\).%
\side{Lies die Gleichung als Satz: \enquote{Wahrscheinlichkeit gleich Spur von Zustand mal Effekt.}
Farbe: \textcolor{tpc-state}{\textbf{Zustand}}, \textcolor{tpc-effect!80!black}{\textbf{Effekt}}.}

\begin{postulat}{P4\enspace Bornsche Regel}
Wird ein System im Zustand \(\St{\rho}\) präpariert und mit dem Effekt \(\Ef{E}\) getestet, so ist die
Wahrscheinlichkeit, dass der Detektor anspricht,
\begin{equation}
  p(\Ef{E}\mid\St{\rho}) \;=\; \tr\bigl(\St{\rho}\,\Ef{E}\bigr)\in[0,1].
\end{equation}
\end{postulat}

\noindent Für ein Teilchen im Kasten der Breite \(L\) im Eigenzustand \(\psi_2\) und einen Detektor, der auf
das Intervall \([a,b]\) anspricht, gilt
\begin{equation}
  p \;=\; \int_a^b \lvert\psi_2(x)\rvert^2\,\dd x
    \;=\; \frac{b-a}{L} - \frac{1}{4\pi}\Bigl[\sin\frac{4\pi x}{L}\Bigr]_a^b
    \;\approx\; \num{0.34}
\end{equation}
für \(a=\num{0.55}\,L\), \(b=\num{0.80}\,L\). Abbildung~1 zeigt dieselben Zahlen.%
\side{\textbf{Live:} \texttt{tpc/6.3-born} \\ Schieberegler für \(a\), \(b\) und \(n\); dieselben Farben, dieselbe Achse.}

\begin{center}
\begin{tikzpicture}[x=78mm,y=9mm,font=\sffamily\scriptsize]
  % forbidden region (environment role): hatched, grey
  \fill[pattern=north east lines,pattern color=tpc-ink-muted!70] (-0.07,0) rectangle (0,2.25);
  \fill[pattern=north east lines,pattern color=tpc-ink-muted!70] (1,0) rectangle (1.07,2.25);
  \draw[tpc-ink,line width=0.6pt] (-0.07,0)--(1.07,0);
  \draw[tpc-ink,line width=0.6pt] (0,0)--(0,2.25) (1,0)--(1,2.25);
  % effect window
  \fill[tpc-effect,opacity=0.35] (0.55,0) rectangle (0.80,2.25);
  \draw[tpc-effect,line width=0.8pt] (0.55,0)--(0.55,2.25) (0.80,0)--(0.80,2.25);
  % state density |psi_2|^2 = 2 sin^2(2 pi x)
  \draw[tpc-state,line width=1.1pt,domain=0:1,samples=121,smooth] plot (\x,{2*sin(360*\x)^2});
  % direct labels
  \node[text=tpc-state,anchor=south west] at (0.27,2.02) {$\lvert\psi_2\rvert^2$};
  \node[text=tpc-effect!75!black,anchor=south] at (0.675,2.25) {Effekt $E$: Detektor auf $[a,b]$};
  \node[text=tpc-ink-muted,anchor=south] at (-0.035,2.25) {$V=\infty$};
  \node[text=tpc-ink-muted,anchor=south] at (1.035,2.25) {$V=\infty$};
  \foreach \t/\l in {0/0,0.55/a,0.8/b,1/L} \node[anchor=north] at (\t,0) {$\l$};
  \node[anchor=north] at (0.5,-0.25) {$x$};
\end{tikzpicture}\\[1mm]
\parbox{108mm}{\sffamily\fontsize{7.8}{10.4}\selectfont\color{tpc-ink-muted}%
\textbf{\color{tpc-ink}Abb.~1\enspace Die schraffierte Fläche ist verboten, die farbige Fläche ist der Detektor.}
Der Anteil der Fläche unter \(\lvert\psi_2\rvert^2\) im Fenster ist \(p\approx\num{0.34}\).}
\end{center}

\begin{zoomin}{6.V1}{T-06.3a}
\textbf{Vertiefung 6.V1 (Warum gerade die Spur?).} Ein Messgerät mit den Ausgängen \(k=1,2,\dots,K\) wird durch Effekte
\(\Ef{E}_k\ge 0\) mit \(\sum_k \Ef{E}_k=\mathbb{1}\) beschrieben (ein POVM). Wegen \(\St{\rho}\ge0\) ist jedes
\(\tr(\St{\rho}\Ef{E}_k)\ge0\), und \(\sum_k\tr(\St{\rho}\Ef{E}_k)=\tr\St{\rho}=1\): Die Bornsche Regel
liefert also für jedes Messgerät eine Wahrscheinlichkeitsverteilung. Umgekehrt ist jede
affine Abbildung von Zuständen nach \([0,1]\) von dieser Form; das ist der Inhalt von P3. \emph{Zurück:} Für den Kasten heißt das nur, dass \(p\) eine Fläche ist.
\end{zoomin}

\noindent Die Zeitentwicklung ändert \(\St{\rho}\), nicht die Regel:
\(\St{\rho}(t)=\eu^{-\iu Ht/\hbar}\,\St{\rho}\,\eu^{\iu Ht/\hbar}\), und \(p(t)=\tr(\St{\rho}(t)\Ef{E})\).
\end{document}
